%!TEX root = ../../main.tex
\section{학습기와 정보 구성에 관하여 말할 수 있는 것}
\label{sec:pred-learner}

M-RQ3의 두 번째 절은 관측 정보의 구성과 학습기의 선택이 예측 성능과 확률 품질에 어떻게 반영되는가이다. 동결 설계가 이 절에 답하는 범위는 제한적이며 그대로 밝힌다.

\paragraph{학습기.} 로지스틱 $q$는 \dimInputQ{}개 입력을, 진단용 LightGBM 후보는 챔피언 ID 10개를 더한 \dimInputLgbm{}개 입력을 썼다. 따라서 둘은 동일 정보의 학습기 비교가 아니다. 본 논문은 로지스틱 사양을 사전에 고정했고 아키텍처 탐색을 하지 않았으므로, 학습기에 관한 순위나 우열을 주장하지 않는다. 선행 연구의 정합 입력 비교(같은 테이블 입력의 그래디언트 부스팅과 MLP)는 다른 라벨과 코퍼스에서의 결과이므로 여기에 옮겨 오지 않는다.

\paragraph{정보 구성.} 학위논문 연구계획은 초기 상태 요약 $C$($\ppre$와 시간 기본 열), $C + F$(프레임 기반 자원·성장·체력 특징군), $C + E$(사건·이력 특징군), 전체 $X$의 네 정보집합(M-F0--M-F3)을 같은 학습기로 비교하는 심화 모듈을 제안하였다. 이 모듈은 실행되지 않았다. 따라서 어떤 특징군이 초기 승률 요약에 조건부로 추가 효용을 갖는지에 관해 본 논문은 결과를 제시하지 않으며, 부록 \ref{app:deferred}에 사전 등록 형식의 계획으로만 둔다. $q$의 계수나 SHAP 같은 특징 귀속을 정보 유무의 독립 검정이나 인과적 승리 조건으로 쓰지 않는다.

\paragraph{확률 품질.} 학습기·정보 구성이 확률 품질에 반영되는 방식은 제 \ref{sec:val-corp} 절의 점수 분해가 $q$와 $\PTflex$ 사이에서 보여 주는 범위까지만 말할 수 있다.
